\section{Tree Models}
\begin{itemize}
	\item Recursively split dataset into subsets.
	\item Each split aims to make the subsets as homogenous as possible wrt. the target.
	\item Preds: Classification $\to$ Majority class in leaf node. Regression $\to$ Mean of target variable of instances in leaf node.
	\item Trees can model nonlinear and feature interactions naturally.
\end{itemize}

\subsection{Decision Trees by CART algorithm}
\textit{Classification and Regression Trees}. Splits determined by Regression: variance of target $y$, Classification: Gini impurity/entropy.
\begin{enumerate}
	\item Test all possible cut-off values for splits
	\item Check variance/impurity reduction per split
	\item Select threshold that achieves the best split
	\item Repeat recursively for child nodes
\end{enumerate}

\paragraph{Prediction}
\begin{equation}
	\hat y = \hat f (x) = \sum_{m=1}^{M}c_m \boldone\lrc{x \in R_m}
	\label{eq:CART-Prediction}
\end{equation}
\begin{itemize}
	\item $c_m$: Prediction in leafe node (Avg $y$ or majority class over training instances in subset $R_m$)
	\item Indicator function: If $x$ is in Region $m$
\end{itemize}
Only one element of the sum will be active.

\paragraph{Splitting Regression}
\begin{equation}
	Var\lr{y} = \frac{1}{N}\displaystyle\sum_{i}(y_i - \bar{y})²
	\label{eq:variance-split}
\end{equation}
\paragraph{Splitting Classificaiton}
\begin{align}
	G = \sum_{k} p_k(1-p_k)   &  & \text{Gini Index} \\
	H = \sum_{k} p_k\log(p_k) &  & \text{Entropy}
\end{align}

\paragraph{Stop Criteria}
We stop splitting, when one or more conditions are met.
\begin{itemize}
	\item Nr of samples in a node before splitting
	\item Nr of saples in a terminal node
	\item Max depth
	\item Not enough reduction in impurity or variance.
\end{itemize}
Final Tree partitions the input space.

\paragraph{Interpreting Decision Trees}
Follow the decision path from root to leaf. Splits along path are connected by logical \texttt{AND}
\begin{itemize}
	\item \textbf{Global Feature importance:} Important features are involved in more splits. And important splits reduce the variance/impurity by a large amount.
	      \begin{equation}
		      \text{Feature importance }= \frac{\text{Total reduction by feature}}{\text{Sum of all reductions}} \cdot 100
		      \label{eq:tree_importance}
	      \end{equation}
	      Where Reduction by feature is the sum of all reductions at each split where the feature is used. Sum of all reductions is the reduction of all features combined.
	\item \textbf{Local:} To get the contribution of a feature to the prediction:
	      \begin{equation}
		      \hat f(x) = \bar y + \sum_{d=1}^{D} \text{split.Contrib}(d,x) = \bar y + \sum_{j=1}^{p} \text{feat.contrib}(j, x)
		      \label{eq:local-tree-contribution}
	      \end{equation}
	      Instance pred =  mean outcome + feature contributions.
	      Contribution could be change in prediction by choosing to take one branch ($|$That subsets mean - supergroups mean$|$).
\end{itemize}

\paragraph{Notes}
\begin{itemize}
	\item Shallow trees (depth $\leq$ 3) provide compact explanations.
	\item Small input changes can cause big jumps in output.
	\item Trees quickly become to difficult to interpret with increasing depth
\end{itemize}

\subsection{Decision Rules}
\begin{itemize}
	\item \texttt{IF} condition \texttt{THEN} prediction.
	\item Rules build from simple statements, can be combined or single rules. Intuitively interpretable
	\item Support: percent of data where rules condition applies
	\item Accuracy: Percentage of instances where the prediction is correct of instances where the rule is applied.
\end{itemize}
Automatically produced by tree ml models

\subsection{RulerFit}
Combines simplicity of linear models with flexibility of decision trees.
Learn a sparse linear model that includes original features and rules capturing nonlinear interactions(Lasso $L1$ norm added).
\paragraph{Rule Generation}
\begin{enumerate}
	\item Fit an ensemble of trees on the data (m trees generates $K = \sum_{m=1}^{M}2(t_m-1)$ rules,
	      where $t_m$ is number of leaf nodes in tree m).
	      Generate trees of different depths for diversity.
	\item Get all unique paths from root to leaf, each one is a decision rule
	\item Use rules as binary features.
	\item Add them as binary features to the model $r = 1$ if rule is active $r=0$
\end{enumerate}
\begin{equation}
	\hat y = \beta_0 + \sum_{j=1}^{p}\beta_jx_j + \sum_{k=1}^{K} \alpha_kr_k
	\label{eq:RulerFit}
\end{equation}







